\section{Fonts and Margins}\label{ch4sec:fonts-margins}

Typography and page geometry are controlled in the class file and the main preamble. This section explains the key variables and how to change them without breaking the layout.

\subsection{Font Selection}

The template uses a standard serif font for body text and a sans-serif font for headings. You can change the font packages in the class file; common choices include \texttt{mathpazo}, \texttt{libertine}, or the default \TeX{} fonts. Table~\ref{ch4tab:demo-arch} lists the main font roles in the template.

\begin{table}[h]
\caption{Font roles in the template}\label{ch4tab:demo-arch}
\centering
\begin{tabular}{l l}
\toprule
\thead{Role} & \thead{Default source} \\
\midrule
Body text & Serif font package in class file \\
Headings & Sans-serif font package in class file \\
Math & \texttt{amsmath}, \texttt{siunitx} setup \\
Code & Monospaced font in listing setup \\
\bottomrule
\end{tabular}
\end{table}

\begin{table}[h]
\caption{Typographic settings}\label{ch4tab:demo-typography}
\centering
\begin{tabular}{l l}
\toprule
\thead{Setting} & \thead{Where to change} \\
\midrule
Font family & Class file font packages \\
Font size & Class file base size options \\
Line spacing & \texttt{setspace} or class file \\
Page margins & \texttt{geometry} options in class file \\
\bottomrule
\end{tabular}
\end{table}

\subsection{Page Margins}

Margins are set using the \texttt{geometry} package. The class file defines reasonable defaults for A4 paper with a binding offset. If your institution requires different margins, update the \texttt{geometry} options in the class file rather than inserting manual page layout commands in individual chapters.
